\section{Related Work}
\paragraph{RL for T2I Models}
% The remarkable success of Reinforcement Learning (RL) in aligning Large Language Models (LLMs) has catalyzed a paradigm shift in visual synthesis. Inspired by these triumphs, recent works have adapted RL frameworks—such as PPO and DPO—to fine-tune Text-to-Image (T2I) models using human-preference reward signals [Cite: DPO-Diffusion, RLHF-V]. While these approaches improve specific metrics like aesthetic quality [Cite: ImageReward, PickScore], they are frequently susceptible to reward hacking, which leads to suboptimal generalization on out-of-distribution (OOD) data.
The success of RL-based alignment in large language models has recently inspired reinforcement learning for text-to-image (T2I) generation. Early methods such as DDPO~\cite{black2024training}, DPOK~\cite{fan2023dpok}, and ImageReward/ReFL~\cite{xu2023imagereward} formulate diffusion generation as policy optimization with rewards for aesthetics, human preference, or text-image alignment, while Diffusion-DPO~\cite{wallace2023diffusiondpo} aligns diffusion models using preference pairs. More recent GRPO-style methods extend RL to modern visual generators, including those for flow models~\cite{liu2025flowgrpo,xue2025dancegrpo}, and AR paradigms~\cite{yuan2025argrpo,zhang2025gcpo,ma2025stage,zhang2025maskfocus,ma2026margrpo} . However, T2I generation requires multiple rewards to cover aesthetics, alignment, fidelity, and compositional correctness. Existing solutions remain hard to control: DanceGRPO~\cite{xue2025dancegrpo} directly mixes rewards such as HPS and CLIP, often trading off one metric against another; Flow-GRPO~\cite{liu2025flowgrpo} uses staged reward/dataset curricula, making results sensitive to ordering and stage design; and GDPO~\cite{liu2026gdpo} shows that GRPO~\cite{guo2025deepseek} may suffer from reward-normalization collapse under multi-reward settings. This motivates a more controllable multi-reward coordination mechanism.

% \paragraph{Multi-Teacher Distillation}
% Multi-teacher distillation (MTD) extends the traditional single-teacher paradigm by leveraging a diverse ensemble of expert models to guide a single student model [Cite: Distilling Knowledge from Multiple Teachers]. The core motivation behind MTD is that a solitary teacher may possess limited or biased knowledge, whereas multiple teachers can provide a more comprehensive and robust supervisory signal across disparate domains. In the field of computer vision and natural language processing, MTD has been successfully applied to enhance model performance by aggregating the logits, intermediate features, or attention maps from various specialized architectures [Cite: Learning from Multiple Teachers with Gradient Harmonization].
% \paragraph{On-Policy Distillation}
% Traditional knowledge distillation (KD) for generative models primarily adheres to an offline framework, where a student model is trained to replicate the output or internal representations of a teacher model using fixed, pre-generated datasets [Cite: Progressive Distillation]. While effective for model compression, offline KD inherently suffers from distributional shift; the teacher's supervision is confined to a static data distribution that may not align with the student's evolving generative trajectories. To mitigate this, On-Policy Distillation (OPD) has emerged as a more robust alternative~\cite{lu2025onpolicydistillation}. By requiring the student to generate samples from its own current policy—which are then evaluated or "corrected" by the teacher—OPD ensures that the supervisory signal remains dynamically coupled with the student's exploration space. This on-policy mechanism significantly improves the fidelity of the distilled model, particularly in high-dimensional tasks where the exploration space is vast.

\paragraph{On-Policy Distillation}
Traditional offline distillation relies on fixed datasets and fails to adapt to the student's evolving trajectory. In contrast, On-Policy Distillation (OPD) dynamically couples the teacher's supervisory signal with the student’s exploration space. In the LLM domain, OPD has seen rapid development: GKD~\cite{agarwal2024policy} established the canonical framework to mitigate exposure bias; MiniLLM~\cite{gu2024minillm} and DistiLLM~\cite{ko2025distillm} introduced Reverse and Skewed KL to refine mode-seeking and optimization stability; G-OPD~\cite{yang2026learning} unified OPD under KL-constrained RL theory; Entropy-Aware OPD~\cite{jin2026entropy} preserves diversity through adaptive divergence functions; Fast OPD~\cite{zhang2026fast} significantly accelerates computation via prefix truncation; and PACED~\cite{xu2026paced} implements a competence-aware curriculum based on gradient signal-to-noise analysis. Despite these LLM advancements, OPD remains underexplored in visual Flow Matching models, which require dense supervision within high-dimensional velocity fields. We propose Flow-OPD, the first systematic migration of on-policy distillation to Flow Matching, utilizing multi-teacher dense supervision to overcome the reward sparsity bottleneck.